\documentclass[12pt,a4paper]{article}

\usepackage[utf8]{inputenc}
\usepackage[T1]{fontenc}
\usepackage{lmodern}
\usepackage[margin=1in]{geometry}
\usepackage{setspace}
\usepackage{microtype}

% --- Abbreviations (keeps the text consistent) ---
\newcommand{\PV}{PV}          % photovoltaic
\newcommand{\PVchar}{P--V}    % power-voltage characteristic

\title{Maximum Power Point Tracking for Photovoltaic Arrays\
under Partial Shading Conditions}
\author{Soumyadip}
\date{}

\begin{document}

\maketitle

\begin{abstract}
\onehalfspacing
\noindent
Reliable and efficient photovoltaic (\PV) power generation is essential for
maximising the utilisation of renewable energy. However, conventional Maximum
Power Point Tracking (MPPT) techniques face significant challenges under
partial shading conditions, where non-uniform irradiance across \PV{} modules
can result in multiple peaks in the power--voltage (\PVchar) characteristic.
This makes it difficult for conventional tracking techniques to consistently
identify and track the Global Maximum Power Point (GMPP), potentially leading
to significant power losses.

This project aims to investigate and develop an effective MPPT strategy for
photovoltaic arrays operating under partial shading conditions. A \PV{} array
integrated with a DC--DC converter will be modelled and simulated to study its
behaviour under both uniform irradiance and different partial shading patterns.
Various MPPT approaches will be investigated and comparatively evaluated, with
conventional techniques such as Perturb and Observe (P\&O) and Incremental
Conductance (INC) studied alongside global optimisation and intelligent
techniques, including Genetic Algorithms (GA), Artificial Neural Networks
(ANNs), and other suitable approaches. The methods will be evaluated based on
parameters such as global maximum power point tracking accuracy, convergence
speed, dynamic response, steady-state oscillations, computational complexity,
and overall energy extraction. Based on the comparative analysis, a suitable
approach will be identified and its effectiveness under different operating
conditions will be assessed. The project ultimately aims to improve the
reliability and efficiency of \PV{} power extraction under partial shading
conditions while establishing an effective balance between tracking performance
and computational requirements.
\end{abstract}

\vspace{1em}
\noindent\textbf{Keywords:} photovoltaic array; maximum power point tracking;
partial shading; global maximum power point; DC--DC converter; perturb and
observe; incremental conductance; genetic algorithm; artificial neural network.

\end{document}
